\documentclass[12pt]{article}
\usepackage[utf8]{inputenc}
\usepackage[T1]{fontenc}
\usepackage{hyperref}
\title{Wikngenchera Language Resources for K-12 Curriculum Development}
\author{Compiled by Claude Sonnet 5.5 (claude-sonnet-5-5) for Andr\'as Kornai}
\date{}
\begin{document}
\maketitle

\section{Introduction}
Wikngenchera (ISO 639-3 wua, Glottocode wikn1246) is a Paman language of the Pama-Nyungan family, spoken on Cape York Peninsula, Queensland, Australia. Search results indicate that it is the same language as Kugu Nganhcara (Kugu-Muminh), which is also called Wik-Nganhcara. Wikipedia reports that only two dialects are still spoken, Kugu-Muminh and Kugu-Uwanh, with about 30 speakers in 2005. The language is classed as severely endangered. Other Wik languages, such as Wik-Mungkan and Wik-Ngathan, are different languages and must not be confused with it.

\section{Primary Educational Resources}
The searches found no primer, textbook, school programme or classroom material for this language. It is not known to be taught in schools.

\section{Linguistic and Reference Materials}
The only sources found are encyclopedic overviews. They give the classification, dialects and the relation to the other Wik languages.
\begin{itemize}
\item Kugu Nganhcara language (Wikipedia): \url{https://en.wikipedia.org/wiki/Kugu_Nganhcara_language}
\item Wik languages (Wikipedia): \url{https://en.wikipedia.org/wiki/Wik_languages}
\end{itemize}
The searches turned up no dictionary or grammar for Kugu Nganhcara under any of its names. A search for them in library and archive catalogues is a sensible next step.

\section{Religious and Community Resources}
No scripture or community organisation specific to this language turned up in the searches.

\section{Digital and Audio Resources}
No audio, video or corpus for Wikngenchera was found. The search also returned a documentation project for the neighbouring Wik-Ngathan and Wik-Ngatharr dialects. It is a different language variety, so it is not listed as a resource for wua.

\section{Curriculum Development Context}
Very little material is available online. A curriculum would have to begin with new documentation work and with archival searches, working with the remaining speakers and the Wik communities.
\end{document}